\section{선별 문제 풀이 II}
\label{sec:galois-gateway-solutions-b}

\subsection*{문제 \thechapter.12: 자동동형 개수 판정}

\begin{solution}
임의의 유한 대수적 확장에 대해
\[
  \Aut_K(L)
  \subseteq\Hom_K(L,\overline K)
\]
이므로
\[
  |\Aut_K(L)|
  \le [L:K]_{\mathrm s}
  \le [L:K].
\]

$L/K$가 갈루아이면 정규성 때문에 모든 $K$-매장이 $L$의 자동동형이고, 분리성 때문에 분리차수가 보통 차수와 같다. 따라서
\[
  |\Aut_K(L)|=[L:K].
\]

반대로 $|\Aut_K(L)|=[L:K]$라 하자. 위 부등식에서 양 끝이 같으므로
\[
  [L:K]_{\mathrm s}=[L:K]
\]
이고 $L/K$는 분리확장이다. 또한 $K$-매장의 총수가 이미 $[L:K]$개인데 그만큼의 자동동형이 존재하므로 모든 $K$-매장은 $L$을 자기 자신으로 보낸다. 체매장 판정에 의해 $L/K$는 정규이다. 따라서 $L/K$는 갈루아 확장이다.
\end{solution}

\subsection*{문제 \thechapter.15: 중간체와 제한 사상}

\begin{solution}
(a) $L/K$가 정규이면 기준체를 확대한 $L/E$도 정규이다. 또한 $L/K$가 분리이면 $L/E$도 분리이다. 따라서 $L/E$는 갈루아 확장이다. $E$를 고정하는 자동동형은 특히 $K$를 고정하므로
\[
  \Gal(L/E)\le\Gal(L/K).
\]

(b) $E/K$도 갈루아라고 하자. $\sigma\in\Gal(L/K)$의 제한 $\sigma|_E$는 $E$의 $K$-매장이다. $E/K$가 정규이므로 $\sigma(E)=E$이고
\[
  \operatorname{res}_E:\Gal(L/K)\to\Gal(E/K)
\]
가 정의된다.

임의의 $\tau\in\Gal(E/K)$는 매장 연장 정리에 의해 $L$의 $K$-매장으로 연장된다. $L/K$가 정규이므로 이 연장은 $L$의 자동동형이다. 따라서 제한 사상은 전사이다. 그 핵은 $E$를 점별로 고정하는 자동동형들의 집합이므로
\[
  \ker(\operatorname{res}_E)=\Gal(L/E).
\]

(c) 군준동형의 기본정리로
\[
  \Gal(L/K)/\Gal(L/E)
  \cong\Gal(E/K)
\]
를 얻는다.
\end{solution}

\subsection*{문제 \thechapter.18: 아르틴의 고정체 정리}

\begin{solution}
$K=L^H$, $n=|H|$라 하자. $a\in L$의 서로 다른 $H$-켤레를 $a_1,\dots,a_r$라 하면
\[
  P_a(X)=\prod_{i=1}^r(X-a_i)
\]
의 근 집합은 $H$에 의해 순열된다. 따라서 모든 계수는 $H$에 의해 고정되고
\[
  P_a\in K[X].
\]
이 다항식은 서로 다른 일차인수의 곱이므로 분리다항식이다. 그러므로 모든 $a\in L$는 $K$ 위에서 분리적이고
\[
  [K(a):K]\le r\le n.
\]

$E/K$가 $L/K$ 안의 유한 부분확장이면 $E/K$는 유한 분리확장이므로 $E=K(b)$인 원시원 $b$가 존재한다. 따라서
\[
  [E:K]=[K(b):K]\le n.
\]
모든 유한 부분확장의 차수는 $n$ 이하이므로 차수가 최대인 유한 부분확장 $E$를 택하자. 임의의 $a\in L$에 대해 $E(a)/K$도 유한 분리확장이고 차수가 $n$ 이하이다. 최대성에서
\[
  [E(a):K]=[E:K]
\]
이므로 $a\in E$이다. 따라서 $L=E$이고 $L/K$는 유한이며
\[
  [L:K]\le n.
\]

정의상 $H\subseteq\Aut_K(L)$이므로
\[
  n=|H|
  \le|\Aut_K(L)|
  \le[L:K]
  \le n.
\]
모든 부등식이 등호가 되어
\[
  [L:K]=n,
  \qquad
  \Aut_K(L)=H.
\]
자동동형 개수 판정에 의해 $L/K$는 갈루아 확장이다.
\end{solution}

\subsection*{문제 \thechapter.19: 유한 갈루아 대응}

\begin{solution}
$G=\Gal(L/K)$라 하자. 두 사상을
\[
  \Phi(H)=L^H,
  \qquad
  \Psi(E)=\Gal(L/E)
\]
로 정의한다.

부분군 $H\le G$에 대해 아르틴의 고정체 정리를 $H$의 $L$ 위 작용에 적용하면
\[
  \Gal(L/L^H)=H.
\]
즉
\[
  (\Psi\circ\Phi)(H)=H.
\]

중간체 $E$에 대해서는 $L/E$가 유한 갈루아 확장이다. 전체 갈루아군의 고정체가 기준체라는 따름정리를 적용하면
\[
  L^{\Gal(L/E)}=E.
\]
즉
\[
  (\Phi\circ\Psi)(E)=E.
\]
따라서 $\Phi$와 $\Psi$는 서로 역인 전단사이다. 고정해야 할 자동동형이 많아질수록 고정체가 작아지므로 대응은 포함관계를 뒤집는다.
\end{solution}

\subsection*{문제 \thechapter.21: 정규부분군과 갈루아 중간체}

\begin{solution}
$E=L^H$라 하자. $E/K$는 $L/K$의 부분확장이므로 분리확장이다. 따라서 갈루아성은 정규성과 동치이다.

먼저 $E/K$가 정규라고 하자. 그러면 모든 $\sigma\in G$가 $E$를 자기 자신으로 보낸다. $\tau\in H$와 $a\in E$에 대해
\[
  (\sigma\tau\sigma^{-1})(a)
  =\sigma\bigl(\tau(\sigma^{-1}(a))\bigr)
  =a.
\]
따라서 $\sigma H\sigma^{-1}\subseteq H$이고 역포함도 같으므로 $H\trianglelefteq G$이다.

반대로 $H\trianglelefteq G$라 하자. $a\in E$, $\sigma\in G$, $\tau\in H$에 대해
\[
  \tau(\sigma(a))
  =\sigma\bigl((\sigma^{-1}\tau\sigma)(a)\bigr)
  =\sigma(a),
\]
왜냐하면 $\sigma^{-1}\tau\sigma\in H$이기 때문이다. 따라서 $\sigma(a)\in E$이고 모든 $\sigma\in G$가 $E$를 보존한다. $E$의 모든 $K$-매장은 $L$의 $K$-자동동형으로 연장되므로 그 상이 $E$이고, 따라서 $E/K$는 정규이다.

이 경우 제한 사상
\[
  G\to\Gal(E/K)
\]
은 전사이고 핵은 $E$를 고정하는 군 $H$이다. 그러므로
\[
  G/H\cong\Gal(E/K).
\]
\end{solution}

\subsection*{문제 \thechapter.23: $S_3$ 확장의 중간체 격자}

\begin{solution}
$G\cong S_3$의 부분군은 자명부분군, 세 개의 위수 $2$ 부분군, $A_3$, 전체 군이다. 갈루아 대응과 차수 공식을 적용하면 다음 표를 얻는다.
\[
\begin{array}{c|c|c|c}
H & L^H & [L^H:\Q] & \Q\text{ 위 정규?}\\
\hline
\{1\}&L&6&\text{예}\\
\langle\tau\rangle&\Q(\alpha)&3&\text{아니오}\\
\sigma\langle\tau\rangle\sigma^{-1}&\Q(\omega\alpha)&3&\text{아니오}\\
\sigma^2\langle\tau\rangle\sigma^{-2}&\Q(\omega^2\alpha)&3&\text{아니오}\\
A_3=\langle\sigma\rangle&\Q(\omega)&2&\text{예}\\
G&\Q&1&\text{예}
\end{array}
\]

예를 들어 $\tau$는 $\alpha$를 고정하므로 $\Q(\alpha)\subseteq L^{\langle\tau\rangle}$이다. 두 체의 $\Q$ 위 차수가 모두 $3$이므로 등호이다. 나머지 두 삼차체는 이를 $\sigma$와 $\sigma^2$로 옮겨 얻는다. $A_3$는 $\omega$를 고정하고 고정체 차수가 $2$이므로 $L^{A_3}=\Q(\omega)$이다.

정규부분군은
\[
  \{1\},\quad A_3,\quad S_3
\]
뿐이다. 따라서 세 삼차체는 비정규이고, 이차체 $\Q(\omega)$는 정규이다. 몫군은
\[
  S_3/A_3\cong C_2
  \cong\Gal(\Q(\omega)/\Q)
\]
이다. 자명한 양 끝에서는 각각 $S_3/\{1\}\cong S_3$와 $S_3/S_3=1$을 얻는다.
\end{solution}

\subsection*{문제 \thechapter.24: 유한체의 갈루아 대응}

\begin{solution}
$G=\Gal(\F_{q^n}/\F_q)=\langle\varphi\rangle$라 하자. 여기서
\[
  \varphi(a)=a^q,
  \qquad |G|=n.
\]
순환군 $C_n$의 모든 부분군은 $d\mid n$에 대해
\[
  H_d=\langle\varphi^d\rangle
\]
꼴이며 그 위수는 $n/d$이다.

$H_d$의 고정원소는 특히 $\varphi^d(a)=a$, 즉
\[
  a^{q^d}=a
\]
를 만족한다. 이 원소들은 정확히 $\F_{q^d}$를 이룬다. 따라서
\[
  \F_{q^n}^{H_d}=\F_{q^d}.
\]
갈루아 대응이 모든 중간체를 부분군에서 얻으므로 중간체는 정확히 $\F_{q^d}$, $d\mid n$이다. 또한
\[
  \Gal(\F_{q^n}/\F_{q^d})=H_d=\langle\Frob_q^d\rangle.
\]
$G$가 아벨군이므로 모든 부분군이 정규이고, 따라서 모든 중간체 $\F_{q^d}/\F_q$가 갈루아이다.
\end{solution}

\subsection*{문제 \thechapter.28: $\Q(\zeta_5)$의 실수 고정체}

\begin{solution}
$\zeta=\zeta_5$라 쓰자. 최소다항식은
\[
  \Phi_5(X)=X^4+X^3+X^2+X+1
\]
이고 차수는 $4$이다. 그 네 근 $\zeta,\zeta^2,\zeta^3,\zeta^4$가 모두 $L=\Q(\zeta)$에 들어 있으므로 $L$은 $\Phi_5$의 분해체이다. 표수 $0$에서 $\Phi_5$는 분리이므로 $L/\Q$는 차수 $4$의 갈루아 확장이다.

복소켤레 $c$는 위수 $2$이므로 갈루아 대응에서
\[
  [L^{\langle c\rangle}:\Q]
  =[\Gal(L/\Q):\langle c\rangle]
  =2.
\]

원소
\[
  u=\zeta+\zeta^{-1}
\]
는 복소켤레에 의해 고정된다. $\Phi_5(\zeta)=0$을 $\zeta^2$로 나누면
\[
  \zeta^2+\zeta+1+\zeta^{-1}+\zeta^{-2}=0.
\]
그런데
\[
  \zeta^2+\zeta^{-2}=u^2-2
\]
이므로
\[
  u^2+u-1=0.
\]
이 이차식의 판별식은 $5$이고 $u\notin\Q$이므로
\[
  [\Q(u):\Q]=2.
\]
따라서 차수 비교로
\[
  L^{\langle c\rangle}=\Q(u).
\]
또한
\[
  u=\frac{-1+\sqrt5}{2}
\]
또는 그 켤레이므로
\[
  \Q(u)=\Q(\sqrt5).
\]
결국
\[
  L^{\langle c\rangle}
  =\Q(\zeta_5+\zeta_5^{-1})
  =\Q(\sqrt5).
\]
\end{solution}
